\documentclass[11pt,a4paper]{article}
% Shared preamble for the English probability and statistics documents.
% Based on the original course preamble. Compile with pdfLaTeX.
\usepackage[top=1.5cm,bottom=2cm,left=2cm,right=2cm]{geometry}
\usepackage[T1]{fontenc}
\usepackage{lmodern}
\usepackage[utf8]{inputenc}
\usepackage[english]{babel}

\usepackage{amsmath}
\usepackage{amssymb}
\usepackage{amsthm}
\usepackage{mathrsfs}
\usepackage{mathtools}
\usepackage{etoolbox}
\usepackage{tikz}
\usetikzlibrary{shapes.multipart,arrows.meta}
\usepackage{xcolor}
\usepackage{enumerate}
\usepackage[shortlabels]{enumitem}
\usepackage{hyperref}
\usepackage{version}
\usepackage{bbold}
\usepackage{fontawesome5}
\usepackage{fancyhdr}
\lfoot{}
\usepackage{natbib}
\usepackage{graphicx}
\usepackage{pdfpages}
\usepackage{authblk}
    
\usepackage{blindtext}
\usepackage{multicol}
\usepackage{collcell}
\usepackage{wrapfig}
\usepackage{blkarray}
\usepackage{cancel}
\usepackage{appendix}
\usepackage{float}
\makeatletter
\input{tex/assets/slashbox.sty}
\makeatother
\hypersetup{unicode,colorlinks,citecolor=red,filecolor=black,
  linkcolor=red,urlcolor=cyan}
\setlength{\columnsep}{1cm}
\pagenumbering{gobble}

\def\numtext#1{%
  \ifcase#1\or one\or two\or three\or four\or five\or six\or seven\or eight\or nine\or ten\or
  eleven\or twelve\or thirteen\or fourteen\or fifteen\or sixteen\or seventeen\or eighteen\or nineteen\or twenty\or Lots\fi}

\newcommand{\Byte}[1]{%
  \def\nodecontents{}%
  \foreach[count=\l] \k in {#1} {\xappto{\nodecontents}{\noexpand\nodepart{\numtext{\l}}\k}}%
  \node[rectangle split,rectangle split parts=8,rectangle split horizontal,draw] at (2,2){\nodecontents};}

\newcommand{\D}{{\mathrm{d}}}
\newcommand{\cE}{{\mathcal{E}}}
\newcommand{\cG}{{\mathcal{G}}}
\newcommand{\cX}{{\mathcal{X}}}
\newcommand{\cY}{{\mathcal{Y}}}
\newcommand{\cB}{{\mathcal{B}}}
\newcommand{\cF}{{\mathcal{F}}}
\newcommand{\cO}{{\mathcal{O}}}
\newcommand{\cS}{{\mathcal{S}}}
\newcommand{\cT}{{\mathcal{T}}}
\newcommand{\given}{\, \big| \,}
\newcommand{\cL}{{\mathscr{L}}}
\newcommand{\cD}{{\mathcal{D}}}
\newcommand{\cA}{{\mathscr{A}}}
\newcommand{\cP}{{\mathscr{P}}}
\newcommand{\cR}{{\mathscr{R}}}
\newcommand{\A}{{\mathcal A}}
\newcommand{\one}{\mathbb{1}}
\newcommand{\mathbbm}[1]{\mathbb{#1}}

\providecommand{\C}{}
\renewcommand{\C}{\mathbb C}
\newcommand{\R}{\mathbb R}
\newcommand{\Q}{\mathbb Q}
\newcommand{\N}{\mathbb N}
\newcommand{\T}{\mathbb T}
\newcommand{\F}{\mathbb F}
\renewcommand{\P}{\mathbb P}
\newcommand{\Z}{\mathbb Z}
\newcommand{\E}{\mathbb {E}}
\newcommand{\Pbb}[1]{\P\left( #1\right)}
\newcommand{\Ebb}[1]{\E\left[ #1\right]}

\newcommand{\lb}{\left (}
\newcommand{\rb}{\right )}
\newcommand{\lbb}{\left [}
\newcommand{\rbb}{\right ]}
\newcommand{\labs}{\left |}
\newcommand{\rabs}{\right |}
\newcommand{\lbrb}[1]{\lb #1 \rb}
\newcommand{\lbbrbb}[1]{\lbb#1\rbb}
\newcommand{\labsrabs}[1]{\labs#1\rabs}
\newcommand{\lbbrb}[1]{\lbb#1\rb}
\newcommand{\lbrbb}[1]{\lb#1\rbb}
\newcommand{\langlerangle}[1]{\langle#1\rangle}
\newcommand{\lbcurly}{\{}
\newcommand{\rbcurly}{\}}
\newcommand{\curly}[1]{\lbcurly#1\rbcurly}
\newcommand{\abs}{\labsrabs}
\newcommand{\ind}[1]{1_{#1}}
\newcommand{\indep}{\perp \!\!\! \perp}
\renewcommand{\d}{\, \mathrm{d}}

\newtheorem{theorem}{Theorem}
\newtheorem{problemCnt}{Lemma}
\theoremstyle{definition}
\newtheorem{problem}[problemCnt]{Problem}
\newtheorem{problembook}[problemCnt]{\faBook \,  Problem}
\newtheorem{problempencil}[problemCnt]{\faPencil* \,  Problem}
\newtheorem{problemstar}[problemCnt]{\faStar\,  Problem}
\newtheorem{problemstarpencil}[problemCnt]{\faStar \, \faPencil* \,   Problem}
\newtheorem{probleminteresting}[problemCnt]{\faSmile[regular] \,   Problem}
\newtheorem{solution}{Solution}

\newcommand{\fwcell}[1]{\makebox[\arraystretch\normalbaselineskip]{$#1$}}
\newenvironment{enumeratei}{\begin{enumerate}[(i)]}{\end{enumerate}}
\newenvironment{enumeratea}{\begin{enumerate}[(a)]}{\end{enumerate}}
\newenvironment{enumeraten}{\begin{enumerate}[\upshape 1.]}{\end{enumerate}}

% Compatibility helpers needed by a few files in the combined archive.
\AtBeginDocument{\noindent{\footnotesize English translation prepared with ChatGPT (OpenAI, GPT-6), 26 October 2026, for translation of the Bulgarian text into English and LaTeX formatting.
The Bulgarian originals remain the reference; translation may contain errors.}\par\smallskip}
\newcommand{\MissingGraphic}[1]{%
  \begin{center}\fbox{\parbox{0.82\linewidth}{\centering Missing source image: \texttt{#1}}}\end{center}}



\pagestyle{fancy}\fancyhf{}\lhead{Probability and Statistics: Resit solutions}\rhead{18 August 2021}
\setlength{\headheight}{15pt}\pagenumbering{arabic}\cfoot{\thepage}
\begin{document}
\noindent\emph{English transcription of the original handwritten solution outline. Its calculations and incomplete final answers are retained; apparent source slips are identified in notes.}

\section*{Problem 1}
Tourists form $2/3$ of the visitors and answer correctly with probability $1/6$; locals form $1/3$ and answer correctly with probability 1. Let $L$ and $R$ denote the events that left and right, respectively, are the correct directions.
By Bayes' formula,
\begin{align*}
\P(L\mid\text{answer }L)
&=\frac{\P(\text{answer }L\mid L)\P(L)}{\P(\text{answer }L)}\\
&=\frac{(\frac23\cdot\frac16+\frac13)\cdot\frac12}
{\P(\text{answer }L\mid L)\P(L)+\P(\text{answer }L\mid R)\P(R)}\\
&=\frac{4/9}{4/9+\frac23\cdot\frac56}=\frac49.
\end{align*}
Similarly,
\begin{align*}
\P(L\mid\text{answers }LL)
&=\frac{\frac23(\frac16)^2+\frac13}
{\frac23(\frac16)^2+\frac13+\frac23(\frac56)^2}\\
&=\frac{2+36}{38+50}=\frac{38}{88}=\frac{19}{44}.
\end{align*}

\section*{Problem 2}
Model bit $i$, $i=1,\ldots,8$, by a Bernoulli random variable $X_i$ with probability $2^{-i}$ of being 1. The expected number of 1-bits is
\[
\E(X_1+\cdots+X_8)=\frac12+\frac14+\cdots+\frac1{2^8}=\frac{255}{256}.
\]
For the expected numerical value,
\begin{align*}
\E[X_1\cdot2^7+X_2\cdot2^6+\cdots+X_8\cdot2^0]
&=\frac12\cdot2^7+\frac14\cdot2^6+\cdots+\frac1{2^8}\cdot1\\
&=2^6+2^4+2^2+2^0+2^{-2}+2^{-4}+2^{-6}+2^{-8}.
\end{align*}
The original's reported numerical answer is $85.328125$. See Note 1.

\section*{Problem 3}
Recall that $f_U(t)=\mathbf1_{[0,1]}(t)$ and $\P(U>0.75)=0.25=:p$.
\subsection*{Method I}
The number $N$ of observations greater than 0.75 is $\operatorname{Bin}(10,p)$, so
\[
\E N=10p=2.5,\qquad \operatorname{Var}(N)=10p(1-p)=1.875.
\]
The estimate obtained by $A$ and $B$ is the proportion $N/10$. Hence
\[
\E(N/10)=\frac{\E N}{10}=0.25,
\qquad \operatorname{Var}(N/10)=\frac{1.875}{100}.
\]
\subsection*{Method II}
Here the estimate is $\widetilde X=\operatorname{Bin}(1000,p)/1000$. Thus
\[
\E\widetilde X=0.25,\qquad
\operatorname{Var}(\widetilde X)=\frac{10^3p(1-p)}{10^6}
=\frac{1.875}{10^4}.
\]
\subsection*{3. Rounds}
\[
\P(\text{both numbers}>0.75)=\left(\frac14\right)^2=\frac1{16}.
\]
By the strong law of large numbers, this happens in a proportion $1/16$ of the rounds.
Also,
\[
\P(\text{at least one number}>0.75)=1-\left(\frac34\right)^2=\frac7{16},
\]
so the expected number of rounds is $\E[\operatorname{Ge}(7/16)]=16/7$.

\section*{Problem 4}
\paragraph{1.} The original writes
\begin{align*}
\P(N_1(-10,10)>N_2(0,100))
&=\P(N(10,110)>0)\\
&=\P\left(N(0,1)>\frac{10}{\sqrt{110}}\right)\approx\cdots.\tag{Note 2}
\end{align*}
\paragraph{2.} The final line of the original outline is
\[
\P(N_1(-10,10)<-0.1).\tag{Note 3}
\]

\section*{Transcription notes}
\begin{enumerate}
\item The displayed eight-term sum is $85.33203125$. The original's $85.328125$ omits the final $2^{-8}$ contribution.
\item The intermediate normal distribution for the difference should have mean $-10$. The source writes $+10$, while its final standard-normal threshold is consistent with mean $-10$.
\item The problem measures the change in percentage points. A reduction from 250,000 to 225,000 euros is $-10$ percentage points, whereas the outline writes the fractional change $-0.1$. The outline does not finish the numerical calculations or provide the second case.
\end{enumerate}
\end{document}
